% =============================================================================
% Bloque 2: Introducción (slides 3–4)
% =============================================================================

\section{Introducción}

% --- Slide 3: Hook + definición + caso real (sin spoiler) ---
\begin{frame}{¿Cuántas muestras se necesitan para clasificar texto?}
\begin{columns}[T]

\begin{column}{0.55\textwidth}
\small
Imagina entrenar un clasificador de \textbf{6 emociones} con apenas\\
\textbf{10 ejemplos por clase}.
\begin{itemize}\setlength\itemsep{3pt}
  \item Total: \alert{60 muestras} en un espacio de \alert{768 dimensiones}.
  \item SMOTE: interpola vectores; no produce texto coherente.
  \item Un LLM genera miles de oraciones plausibles\dots\\
        \cita{pero, ¿son útiles todas?}
\end{itemize}

\vspace{0.5em}
\begin{insight}
\scriptsize
Aprender fronteras de decisión con $60$ puntos en $\mathbb{R}^{768}$ es estadísticamente débil. La aumentación de datos no es opcional: es \textbf{necesaria}.
\end{insight}
\end{column}

\begin{column}{0.45\textwidth}
\setbeamercolor{block title}{fg=white,bg=uaibluedark}
\begin{block}{\small Few-shot text classification}
\scriptsize
Aprender a clasificar texto con un número reducido de ejemplos etiquetados por clase, típicamente $k \leq 50$.\\[3pt]
\textbf{Esta tesis estudia $k \in \{10, 25, 50\}$.}
\end{block}

\vspace{0.6em}
\setbeamercolor{block title}{fg=white,bg=uaiorange}
\begin{block}{\small Caso real: \dataset{emotion}}
\scriptsize
6 clases: alegría, tristeza, miedo, ira, amor, sorpresa.\\
Entrenamiento: \alert{\textbf{60 mensajes}}. Evaluación: $\sim$1.000.\\[2pt]
\textit{La situación del ejemplo es real y se mide en esta tesis.}
\end{block}

\vspace{0.5em}
\centering
{\scriptsize\itshape ¿Cómo aumentar sin etiquetar miles más?}
\end{column}

\end{columns}
\end{frame}

% --- Slide 4: Escena situacional (donde esta tesis es útil) ---
\begin{frame}{Pocos datos, apuro y cero presupuesto: una escena habitual}
\begin{columns}[T]

\begin{column}{0.55\textwidth}
\small
\textbf{Lunes, 9:00.} Tu jefe te pide un clasificador en producción \textbf{el viernes}.

\vspace{0.5em}
\textbf{Lo que tienes:}
\begin{itemize}\setlength\itemsep{2pt}
  \item 12 categorías de mensajes de pacientes (salud digital).
  \item $\sim$15 mensajes etiquetados por categoría desde el CRM.
  \item Cero presupuesto para etiquetado externo.
  \item Un experto clínico ocupado hasta fin de mes.
\end{itemize}

\vspace{0.4em}
\textbf{Lo que el cliente exige:}
\begin{itemize}\setlength\itemsep{2pt}
  \item Macro-F1 $\geq 75\%$ sobre las 12 categorías.
  \item Demo funcionando en \textbf{4 días}.
\end{itemize}

\vspace{0.4em}
\cita{La escena se repite: spam emergente, soporte en idiomas nuevos, triaje clínico\dots}
\end{column}

\begin{column}{0.45\textwidth}
\centering
\vspace{0.4em}
\begin{tikzpicture}[every node/.style={font=\scriptsize, align=center}]
  \node[draw=uaibluedark, fill=white, rounded corners, very thick,
        text width=3.6cm, minimum height=1.1cm] at (0, 3.2)
        {\textbf{Disponible}\\\scriptsize 12 clases $\times$ 15 ej.\\$= 180$ muestras};

  \draw[->, very thick, uaigray] (0, 2.55) -- (0, 2.05);

  \node[draw=uaiorange, fill=uaiorange!10, rounded corners, very thick,
        text width=3.6cm, minimum height=1.1cm] at (0, 1.3)
        {\textbf{Necesario}\\\scriptsize macro-F1 $\geq 75\%$\\en 4 días};

  \node[font=\Large\bfseries, color=uaiorange!85!black] at (0, 0.0) {?};
  \node[font=\small\bfseries, color=uaiorange!85!black] at (0, -0.5) {¿Cómo cerrar la brecha?};
\end{tikzpicture}
\end{column}

\end{columns}

\mensaje{Esta tesis ofrece una respuesta concreta para escenarios como este.}
\end{frame}
